\documentclass[11pt,a4paper]{article}
\usepackage{fontspec,amsmath,amssymb,amsthm,microtype}
\setmainfont{lmroman10-regular.otf}[BoldFont=lmroman10-bold.otf,ItalicFont=lmroman10-italic.otf,BoldItalicFont=lmroman10-bolditalic.otf]
\usepackage{xcolor}
\usepackage[most]{tcolorbox}
\usepackage{mathrsfs}
\providecommand{\lt}{<}
\providecommand{\gt}{>}
\setmonofont{DejaVuSansMono.ttf}[BoldFont=DejaVuSansMono-Bold.ttf,ItalicFont=DejaVuSansMono-Oblique.ttf,BoldItalicFont=DejaVuSansMono-BoldOblique.ttf,Scale=0.86]
\newfontfamily\leanfallback{STIXTwoMath-Regular.otf}[Scale=0.86]
\definecolor{leanrule}{gray}{0.62}
\definecolor{leanback}{gray}{0.962}
\definecolor{leantitle}{gray}{0.22}
\definecolor{leanmuted}{gray}{0.38}
\newtcolorbox{leanbox}[1]{enhanced,breakable,colback=leanback,colframe=leanrule,
  boxrule=0.4pt,arc=0.6mm,left=2.4mm,right=2.4mm,top=1.8mm,bottom=1.8mm,
  before skip=0.8em,after skip=1em,fonttitle=\small\bfseries,coltitle=leantitle,
  colbacktitle=leanback,titlerule=0.3pt,title={#1}}
\newlength{\leancharwidth}
\newenvironment{leancode}{\par\smallskip\begingroup\ttfamily\footnotesize\raggedright
  \settowidth{\leancharwidth}{M}\setlength{\parskip}{0pt}\setlength{\parindent}{0pt}
  \frenchspacing}{\par\endgroup\smallskip}
\newcommand{\leanline}[2]{\hangindent=\dimexpr#1\leancharwidth+4\leancharwidth\relax
  \hangafter=1\noindent\hspace*{#1\leancharwidth}#2\par}
\newcommand{\leansource}[1]{\par{\footnotesize\color{leanmuted}#1}\par}
\newcommand{\leanlabel}[1]{\par\smallskip\noindent{\small\bfseries #1}\par}
\usepackage[margin=25mm]{geometry}
\usepackage[unicode,hidelinks]{hyperref}
\hypersetup{pdfauthor={LangevinLean Project}}
\setlength{\emergencystretch}{3em}
\setlength{\parskip}{0.4em}
\setlength{\parindent}{0pt}
\allowdisplaybreaks[2]
\raggedbottom
\title{\texorpdfstring{MYIPLA: laws, moments and the error bound about the smoothed marginal maximizer}{MYIPLA: laws, moments and the error bound about the smoothed marginal maximizer}}
\author{LangevinLean Project}
\date{}
\begin{document}\maketitle
Draft mathematical exposition; these notes may contain mistakes.


Let \(g_1\) be a differentiable, \(m\)-strongly convex pair potential with \(L_1\)-Lipschitz gradient, let \(g_2\) be convex and bounded below, and let \(\eta>0\). The smoothed pair potential is \(V_\eta=g_1+g_2^\eta\), where \(g_2^\eta\) is the Moreau envelope, and \(L_\eta=L_1+\eta^{-1}\). By the previous sections \(V_\eta\) is differentiable and \(m\)-strongly convex, \(\nabla V_\eta\) is \(L_\eta\)-Lipschitz, and MYIPLA with the proximal map \(P_\eta\) of \(g_2\) is, iterate by iterate, the IPLA scheme for \(V_\eta\). This section turns these identities into statements about the laws \(\rho_n\) of the MYIPLA iterates on a probability space \((\Omega,Q)\).




We prove: \(\rho_n\) is the law of the IPLA iterate for \(V_\eta\); \(\rho_n\) is a probability measure whose parameter block has finite second moment for every \(L^2\) initialization and \(L^2\) innovations; and, for independent standard Gaussian innovations independent of the initialization and \(0<\gamma\le2/(m+L_\eta)\),


\[\Bigl(\mathbb E\|\theta_n-\theta^\star_\eta\|^2\Bigr)^{1/2}\le\sqrt{\frac{d_\theta}{Nm}}+(1-m\gamma)^nW_2(\rho_0,\pi_{N,V_\eta})+\frac{33}{20}\,\frac{L_\eta}{m}\sqrt{d_x+\frac{d_\theta}N}\,\sqrt\gamma,\]


where \(\theta_n\) is the parameter block of the \(n\)-th iterate, \(\theta^\star_\eta\) a maximizer of the smoothed marginal likelihood \(\theta\mapsto\int e^{-V_\eta(\theta,x)}\,dx\) and \(\pi_{N,V_\eta}\) the Gibbs target of the smoothed model. The next section replaces \(\theta^\star_\eta\) by the maximizer of the original nonsmooth model.


\section{\texorpdfstring{Curvature and gradient regularity of the smoothed potential}{Curvature and gradient regularity of the smoothed potential}}


\textbf{Definitions.} For pair potentials \(g_1,g_2\) and \(\eta>0\), the smoothed potential is \(V_\eta=g_1+g_2^\eta\), and for \(L_1\ge0\) the smoothed gradient constant is \(L_\eta=L_1+\eta^{-1}\).




\textbf{Claim.} Let \(E\) be a real normed space (in (b), a complete real inner-product space).



\begin{enumerate}


\item If \(g_1:E\to\mathbb R\) is \(m\)-strongly convex and \(e:E\to\mathbb R\) is convex, then \(g_1+e\) is \(m\)-strongly convex.

\item If \(g_1\) and \(e\) are differentiable, \(\nabla g_1\) is \(L_1\)-Lipschitz and \(\nabla e\) is \(L_2\)-Lipschitz, then \(\nabla(g_1+e)\) is \((L_1+L_2)\)-Lipschitz.

\end{enumerate}



In particular, if \(g_1\) is differentiable, \(m\)-strongly convex with \(L_1\)-Lipschitz gradient, and \(g_2\) is convex and bounded below on a finite-dimensional space, then \(V_\eta\) is \(m\)-strongly convex with \(L_\eta\)-Lipschitz gradient.




\textbf{Proof.} (a) Let \(x,y\in E\), \(a,b\ge0\), \(a+b=1\). Strong convexity of \(g_1\) and convexity of \(e\) give


\[g_1(ax+by)\le ag_1(x)+bg_1(y)-ab\frac m2\|x-y\|^2,\qquad e(ax+by)\le ae(x)+be(y).\]


Adding the two inequalities gives the \(m\)-strong convexity inequality for \(g_1+e\).




(b) By additivity of the gradient of differentiable functions (previous section), \(\nabla(g_1+e)=\nabla g_1+\nabla e\). By the triangle inequality,


\[\|\nabla(g_1+e)(x)-\nabla(g_1+e)(y)\|\le\|\nabla g_1(x)-\nabla g_1(y)\|+\|\nabla e(x)-\nabla e(y)\|\le(L_1+L_2)\|x-y\|.\]


For the final statement apply (a) and (b) with \(e=g_2^\eta\), which is convex and differentiable with \(\eta^{-1}\)-Lipschitz gradient by the Moreau envelope section, so \(L_2=\eta^{-1}\). \(\square\)

\begin{leanbox}{Lean: \texttt{smoothedPotential}, \texttt{smoothingLipschitz}, \texttt{smoothingLipschitz\_\allowbreak{}coe} and 2 more}
\leansource{\texttt{smoothedPotential} \quad ProximalLangevin/\allowbreak{}MYIPLA.lean, lines 24–26}
\begin{leancode}
\leanline{0}{def\ smoothedPotential\ \{E\ :\ Type*\}\ [NormedAddCommGroup\ E]\ [InnerProductSpace\ ℝ\ E]}
\leanline{4}{[FiniteDimensional\ ℝ\ E]\ (g₁\ g₂\ :\ E\ →\ ℝ)\ (η\ :\ ℝ)\ (x\ :\ E)\ :\ ℝ\ :=}
\leanline{2}{g₁\ x\ +\ moreau\ g₂\ η\ x}
\end{leancode}
\leansource{\texttt{smoothingLipschitz} \quad ProximalLangevin/\allowbreak{}MYIPLA.lean, lines 29–29}
\begin{leancode}
\leanline{0}{def\ smoothingLipschitz\ (L₁\ :\ ℝ≥0)\ (η\ :\ ℝ)\ :\ ℝ≥0\ :=\ L₁\ +\ η.toNNReal⁻¹}
\end{leancode}
\leansource{\texttt{smoothingLipschitz\_\allowbreak{}coe} \quad ProximalLangevin/\allowbreak{}MYIPLA.lean, lines 31–33}
\begin{leancode}
\leanline{0}{theorem\ smoothingLipschitz\_coe\ (L₁\ :\ ℝ≥0)\ \{η\ :\ ℝ\}\ (hη\ :\ 0\ <\ η)\ :}
\leanline{4}{(smoothingLipschitz\ L₁\ η\ :\ ℝ)\ =\ (L₁\ :\ ℝ)\ +\ η⁻¹}
\end{leancode}
\leansource{\texttt{smoothed\_\allowbreak{}strongConvex} \quad ProximalLangevin/\allowbreak{}MYIPLA.lean, lines 36–45}
\begin{leancode}
\leanline{0}{theorem\ smoothed\_strongConvex\ \{E\ :\ Type*\}\ [NormedAddCommGroup\ E]}
\leanline{4}{[NormedSpace\ ℝ\ E]\ \{g₁\ e\ :\ E\ →\ ℝ\}\ \{m\ :\ ℝ\}}
\leanline{4}{(hsc\ :\ StrongConvexOn\ univ\ m\ g₁)\ (he\ :\ ConvexOn\ ℝ\ univ\ e)\ :}
\leanline{4}{StrongConvexOn\ univ\ m\ (fun\ x\ ↦\ g₁\ x\ +\ e\ x)}
\end{leancode}
\leansource{\texttt{smoothed\_\allowbreak{}gradient\_\allowbreak{}lipschitz} \quad ProximalLangevin/\allowbreak{}MYIPLA.lean, lines 48–57}
\begin{leancode}
\leanline{0}{theorem\ smoothed\_gradient\_lipschitz\ \{E\ :\ Type*\}\ [NormedAddCommGroup\ E]}
\leanline{4}{[InnerProductSpace\ ℝ\ E]\ [CompleteSpace\ E]}
\leanline{4}{\{g₁\ e\ :\ E\ →\ ℝ\}\ \{L₁\ L₂\ :\ ℝ≥0\}}
\leanline{4}{(hg₁\ :\ Differentiable\ ℝ\ g₁)\ (he\ :\ Differentiable\ ℝ\ e)}
\leanline{4}{(h₁\ :\ LipschitzWith\ L₁\ (gradient\ g₁))\ (h₂\ :\ LipschitzWith\ L₂\ (gradient\ e))\ :}
\leanline{4}{LipschitzWith\ (L₁\ +\ L₂)\ (gradient\ (fun\ x\ ↦\ g₁\ x\ +\ e\ x))}
\end{leancode}
\leanlabel{How the Lean proof follows the argument}
\texttt{smoothed\allowbreak{}Potential g\ensuremath{{}_{1}} g\ensuremath{{}_{2}} \ensuremath{\eta} x} is \texttt{g\ensuremath{{}_{1}} x + moreau g\ensuremath{{}_{2}} \ensuremath{\eta} x}. \texttt{smoothing\allowbreak{}Lipschitz L\ensuremath{{}_{1}} \ensuremath{\eta}} is the nonnegative real \texttt{L\ensuremath{{}_{1}} + \ensuremath{\eta}.\allowbreak{}to\allowbreak{}NNReal\ensuremath{{}^{-1}}}, where \texttt{\ensuremath{\eta}.\allowbreak{}to\allowbreak{}NNReal} is \(\max(\eta,0)\) as an element of \texttt{\ensuremath{\mathbb R}\ensuremath{\ge}0} and inversion in \texttt{\ensuremath{\mathbb R}\ensuremath{\ge}0} sends \(0\) to \(0\); \texttt{smoothing\allowbreak{}Lipschitz\_\allowbreak{}coe} states that for \texttt{h\ensuremath{\eta} : 0 < \ensuremath{\eta}} its real value is \(L_1+\eta^{-1}\) (by \texttt{Real.\allowbreak{}to\allowbreak{}NNReal\_\allowbreak{}of\_\allowbreak{}nonneg}). Using a nonnegative real is required because \texttt{Lipschitz\allowbreak{}With} takes its constant in \texttt{\ensuremath{\mathbb R}\ensuremath{\ge}0}.




\texttt{smoothed\_\allowbreak{}strong\allowbreak{}Convex} is (a). After \texttt{refine \ensuremath{\langle}convex\_\allowbreak{}univ, ?\_\allowbreak{}\ensuremath{\rangle}} and introduction of the points and weights, \texttt{h\ensuremath{{}_{1}}} is the strong convexity inequality \texttt{hsc.\allowbreak{}2} and \texttt{h\ensuremath{{}_{2}}} the convexity inequality \texttt{he.\allowbreak{}2}, and \texttt{linarith} adds them. No sign condition on \(m\) is needed. \texttt{smoothed\_\allowbreak{}gradient\_\allowbreak{}lipschitz} is (b): \texttt{hsum} is \(\nabla(g_1+e)=\nabla g_1+\nabla e\), obtained from \texttt{gradient\_\allowbreak{}add\_\allowbreak{}of\_\allowbreak{}differentiable} by \texttt{funext}, and \texttt{h\ensuremath{{}_{1}}.\allowbreak{}add h\ensuremath{{}_{2}}} is the Mathlib fact that the sum of an \(L_1\)-Lipschitz and an \(L_2\)-Lipschitz map is \((L_1+L_2)\)-Lipschitz.
\end{leanbox}

\section{\texorpdfstring{The MYIPLA laws are IPLA laws and have finite second moments}{The MYIPLA laws are IPLA laws and have finite second moments}}


\textbf{Setting.} Let \((\Omega,Q)\) be a probability space, \(Z_0:\Omega\to\mathbb R^{d_\theta}\times(\mathbb R^{d_x})^N\) an initial state and \(\xi_0,\xi_1,\dots\) innovations, all square-integrable random vectors. For \(\omega\in\Omega\) let \(Z_n(\omega)\) be the MYIPLA chain started at \(Z_0(\omega)\) and driven by \((\xi_k(\omega))_k\). The \emph{law} of the \(n\)-th iterate is the pushforward measure \(\rho_n=Q\circ Z_n^{-1}\), \(\rho_n(A)=Q(Z_n\in A)\). Write \(\theta(z)\) for the parameter block of a state \(z\).




\textbf{Claim.} Let \(N\ge1\), let \(g_1,e\) be differentiable pair functions, \(P\) a map of the pair space and \(\eta\ne0\) with \(\nabla e=\eta^{-1}(I-P)\). Let \(\gamma\in\mathbb R\), \(L\ge0\), and assume that \(\nabla(g_1+e)\) is \(L\)-Lipschitz. Then for every \(n\):



\begin{enumerate}


\item \(\rho_n\) equals the law of the \(n\)-th IPLA iterate for the potential \(g_1+e\) with step \(\gamma\), started at \(Z_0\) and driven by the same innovations;

\item for every \(\vartheta\in\mathbb R^{d_\theta}\), \(\int\|\theta(z)-\vartheta\|^2\,\rho_n(dz)<\infty\).

\end{enumerate}



\textbf{Proof.} (a) By the last identity of the previous section, \(Z_n(\omega)\) coincides for every \(\omega\) with the IPLA iterate \(Y_n(\omega)\) for \(g_1+e\) driven by the same inputs. Equal maps have equal pushforwards, so \(\rho_n=Q\circ Y_n^{-1}\).




(b) Since \(\nabla(g_1+e)\) is \(L\)-Lipschitz, the lifted gradient \(\nabla F_{g_1+e}\) is Lipschitz on the state space, hence of at most linear growth: \(\|\nabla F_{g_1+e}(z)\|\le\|\nabla F_{g_1+e}(0)\|+NL\|z\|\). The IPLA step \(z\mapsto z-\frac\gamma N\nabla F_{g_1+e}(z)+\sqrt{2\gamma/N}\,\xi\) therefore maps square-integrable random vectors to square-integrable random vectors, and by induction on \(n\) the iterate \(Y_n\) is square-integrable (and measurable). The parameter projection \(z\mapsto\theta(z)\) is linear with norm at most \(1\), so \(\theta(Y_n)-\vartheta\) is square-integrable on the probability space \((\Omega,Q)\), constants being square-integrable there. By the change-of-variables formula for pushforward measures,


\[\int\|\theta(z)-\vartheta\|^2\,\rho_n(dz)=\int_\Omega\|\theta(Y_n(\omega))-\vartheta\|^2\,Q(d\omega)<\infty .\ \square\]


No Gaussianity, independence or convexity is used here; only the Lipschitz gradient and the square-integrability of the inputs.

\begin{leanbox}{Lean: \texttt{myiplaLaws}, \texttt{myiplaLaws\_\allowbreak{}eq\_\allowbreak{}iplaLaws}, \texttt{myiplaLaws\_\allowbreak{}parameter\_\allowbreak{}memLp}}
\leansource{\texttt{myiplaLaws} \quad ProximalLangevin/\allowbreak{}MYIPLA.lean, lines 60–65}
\begin{leancode}
\leanline{0}{def\ myiplaLaws\ \{Ω\ :\ Type*\}\ [MeasurableSpace\ Ω]\ (Q\ :\ Measure\ Ω)}
\leanline{4}{\{N\ dθ\ dx\ :\ ℕ\}\ (g₁\ :\ IPLA.Joint\ dθ\ dx\ →\ ℝ)}
\leanline{4}{(P\ :\ IPLA.Joint\ dθ\ dx\ →\ IPLA.Joint\ dθ\ dx)\ (η\ γ\ :\ ℝ)}
\leanline{4}{(initial\ :\ Ω\ →\ IPLA.State\ N\ dθ\ dx)\ (noise\ :\ ℕ\ →\ Ω\ →\ IPLA.State\ N\ dθ\ dx)}
\leanline{4}{(n\ :\ ℕ)\ :\ Measure\ (IPLA.State\ N\ dθ\ dx)\ :=}
\leanline{2}{Q.map\ (fun\ ω\ ↦\ myiplaChain\ g₁\ P\ η\ γ\ (initial\ ω)\ (fun\ k\ ↦\ noise\ k\ ω)\ n)}
\end{leancode}
\leansource{\texttt{myiplaLaws\_\allowbreak{}eq\_\allowbreak{}iplaLaws} \quad ProximalLangevin/\allowbreak{}MYIPLA.lean, lines 68–83}
\begin{leancode}
\leanline{0}{theorem\ myiplaLaws\_eq\_iplaLaws\ \{Ω\ :\ Type*\}\ [MeasurableSpace\ Ω]}
\leanline{4}{(Q\ :\ Measure\ Ω)\ [IsProbabilityMeasure\ Q]\ \{N\ dθ\ dx\ :\ ℕ\}\ (hN\ :\ 0\ <\ N)}
\leanline{4}{\{g₁\ e\ :\ IPLA.Joint\ dθ\ dx\ →\ ℝ\}\ (hg₁\ :\ Differentiable\ ℝ\ g₁)}
\leanline{4}{(he\ :\ Differentiable\ ℝ\ e)\ (P\ :\ IPLA.Joint\ dθ\ dx\ →\ IPLA.Joint\ dθ\ dx)}
\leanline{4}{\{η\ :\ ℝ\}\ (hη\ :\ η\ ≠\ 0)}
\leanline{4}{(hgrad\ :\ ∀\ x,\ gradient\ e\ x\ =\ η⁻¹\ •\ (x\ -\ P\ x))\ (γ\ :\ ℝ)\ (L\ :\ ℝ≥0)}
\leanline{4}{(hLip\ :\ LipschitzWith\ L\ (gradient\ (fun\ x\ ↦\ g₁\ x\ +\ e\ x)))}
\leanline{4}{(initial\ :\ Lp\ (IPLA.State\ N\ dθ\ dx)\ 2\ Q)}
\leanline{4}{(noise\ :\ ℕ\ →\ Lp\ (IPLA.State\ N\ dθ\ dx)\ 2\ Q)\ (n\ :\ ℕ)\ :}
\leanline{4}{myiplaLaws\ Q\ g₁\ P\ η\ γ\ initial\ (fun\ k\ ↦\ noise\ k)\ n\ =}
\leanline{6}{IPLA.laws\ Q\ (fun\ x\ ↦\ g₁\ x\ +\ e\ x)\ γ\ L\ (hg₁.add\ he)\ hLip\ initial\ noise\ n}
\end{leancode}
\leansource{\texttt{myiplaLaws\_\allowbreak{}parameter\_\allowbreak{}memLp} \quad ProximalLangevin/\allowbreak{}MYIPLA.lean, lines 88–106}
\begin{leancode}
\leanline{0}{theorem\ myiplaLaws\_parameter\_memLp\ \{Ω\ :\ Type*\}\ [MeasurableSpace\ Ω]}
\leanline{4}{(Q\ :\ Measure\ Ω)\ [IsProbabilityMeasure\ Q]\ \{N\ dθ\ dx\ :\ ℕ\}\ (hN\ :\ 0\ <\ N)}
\leanline{4}{\{g₁\ e\ :\ IPLA.Joint\ dθ\ dx\ →\ ℝ\}\ (hg₁\ :\ Differentiable\ ℝ\ g₁)}
\leanline{4}{(he\ :\ Differentiable\ ℝ\ e)\ (P\ :\ IPLA.Joint\ dθ\ dx\ →\ IPLA.Joint\ dθ\ dx)}
\leanline{4}{\{η\ :\ ℝ\}\ (hη\ :\ η\ ≠\ 0)}
\leanline{4}{(hgrad\ :\ ∀\ x,\ gradient\ e\ x\ =\ η⁻¹\ •\ (x\ -\ P\ x))\ (γ\ :\ ℝ)\ (L\ :\ ℝ≥0)}
\leanline{4}{(hLip\ :\ LipschitzWith\ L\ (gradient\ (fun\ x\ ↦\ g₁\ x\ +\ e\ x)))}
\leanline{4}{(initial\ :\ Lp\ (IPLA.State\ N\ dθ\ dx)\ 2\ Q)}
\leanline{4}{(noise\ :\ ℕ\ →\ Lp\ (IPLA.State\ N\ dθ\ dx)\ 2\ Q)\ (n\ :\ ℕ)}
\leanline{4}{(θ\ :\ EuclideanSpace\ ℝ\ (Fin\ dθ))\ :}
\leanline{4}{MemLp\ (fun\ z\ ↦\ IPLA.parameterCLM\ z\ -\ θ)\ 2}
\leanline{6}{(myiplaLaws\ Q\ g₁\ P\ η\ γ\ initial\ (fun\ k\ ↦\ noise\ k)\ n)}
\end{leancode}
\leanlabel{How the Lean proof follows the argument}
\texttt{myipla\allowbreak{}Laws Q g\ensuremath{{}_{1}} P \ensuremath{\eta} \ensuremath{\gamma} initial noise n} is \texttt{Q.\allowbreak{}map (fun \ensuremath{\omega} \ensuremath{\mapsto} myipla\allowbreak{}Chain g\ensuremath{{}_{1}} P \ensuremath{\eta} \ensuremath{\gamma} (initial \ensuremath{\omega}) (fun k \ensuremath{\mapsto} noise k \ensuremath{\omega}) n)}; \texttt{Q.\allowbreak{}map f} is the pushforward measure \(Q\circ f^{-1}\). In the two theorems the initialization and innovations are elements of \texttt{Lp (IPLA.\allowbreak{}State N d\ensuremath{\theta} dx) 2 Q}, the space \(L^2(\Omega,Q)\) of square-integrable random states (equivalence classes up to \(Q\)-null sets, applied to \(\omega\) through a chosen representative).




\texttt{myipla\allowbreak{}Laws\_\allowbreak{}eq\_\allowbreak{}ipla\allowbreak{}Laws} is (a). The right-hand side \texttt{IPLA.\allowbreak{}laws Q (fun x \ensuremath{\mapsto} g\ensuremath{{}_{1}} x + e x) \ensuremath{\gamma} L … initial noise n} is the law of the IPLA chain computed in \(L^2\) (the recursion \texttt{IPLA.\allowbreak{}lp\allowbreak{}Chain}); \texttt{IPLA.\allowbreak{}laws\_\allowbreak{}eq\_\allowbreak{}scaled\allowbreak{}Chain} rewrites it as the pushforward of the pointwise chain \texttt{IPLA.\allowbreak{}scaled\allowbreak{}Chain}, after which \texttt{congr 1} and \texttt{funext \ensuremath{\omega}} reduce the goal to \texttt{myipla\allowbreak{}Chain\_\allowbreak{}eq\_\allowbreak{}scaled\allowbreak{}Chain} at each \(\omega\).




\texttt{myipla\allowbreak{}Laws\_\allowbreak{}parameter\_\allowbreak{}mem\allowbreak{}Lp} is (b): \texttt{Mem\allowbreak{}Lp f 2 \ensuremath{\mu}} means that \(f\) is almost-everywhere strongly measurable and \(\int\|f\|^2\,d\mu<\infty\), and \texttt{IPLA.\allowbreak{}parameter\allowbreak{}CLM} is the parameter projection as a continuous linear map. After rewriting with (a), \texttt{chain} is the \(L^2\) iterate; \texttt{mem\allowbreak{}Lp\_\allowbreak{}map\_\allowbreak{}measure\_\allowbreak{}iff} transfers the claim from the pushforward measure to the composition with \texttt{chain} (this is the change-of-variables step), and \texttt{(Lp.\allowbreak{}mem\allowbreak{}Lp chain).\allowbreak{}continuous\allowbreak{}Linear\allowbreak{}Map\_\allowbreak{}comp IPLA.\allowbreak{}parameter\allowbreak{}CLM} followed by \texttt{.\allowbreak{}sub (mem\allowbreak{}Lp\_\allowbreak{}const \ensuremath{\theta})} proves square-integrability of \(\theta(Y_n)-\vartheta\). The linear-growth argument of the text is carried out in the construction of \texttt{IPLA.\allowbreak{}lp\allowbreak{}Chain}, where each step is shown to map \(L^2\) to \(L^2\).
\end{leanbox}


Earlier results used here: \cite[Section~3.1.1]{LangevinLean2026IPLA}; \cite[Section~3.1.4]{LangevinLean2026IPLA}; \cite[Section~3.1.6]{LangevinLean2026IPLA}; \cite[Section~3.3.3]{LangevinLean2026IPLA}; \cite[Section~3.4.3]{LangevinLean2026IPLA}.

\section{\texorpdfstring{The error bound for a general smoothing with residual gradient}{The error bound for a general smoothing with residual gradient}}


\textbf{Definitions.} The \emph{Gibbs target} of a pair potential \(V\) is the probability measure on the state space with density proportional to \(e^{-F_V(z)}\), \(F_V(z)=\sum_iV(J_iz)\); it is denoted \(\pi_{N,V}\). A \emph{marginal maximizer} of \(V\) is a point \(\theta^\star\in\mathbb R^{d_\theta}\) with \(k_V(\theta)\le k_V(\theta^\star)\) for every \(\theta\), where \(k_V(\theta)=\int_{\mathbb R^{d_x}}e^{-V(\theta,x)}\,dx\in[0,\infty]\). For probability measures \(\mu,\nu\) on the state space with finite second moments, \(W_2(\mu,\nu)=\inf_\pi\bigl(\int\|u-v\|^2\,\pi(du,dv)\bigr)^{1/2}\), the infimum over all couplings \(\pi\) of \(\mu\) and \(\nu\). The innovations are \emph{independent standard Gaussian, independent of the initialization} if every coordinate \(\xi_{k,a}\) has the law \(N(0,1)\), the family \((\xi_{k,a})_{k\ge0,\,a}\) of all coordinates is mutually independent, and \(Z_0\) is independent of the sequence \((\xi_k)_{k\ge0}\).




\textbf{Claim.} Let \(N\ge1\), \(m>0\), \(\eta>0\), \(\gamma>0\) and \(L_1\ge0\). Let \(g_1\) be differentiable and \(m\)-strongly convex with \(L_1\)-Lipschitz gradient; let \(e\) be differentiable and convex with \(\eta^{-1}\)-Lipschitz gradient, and let \(P\) be a map of the pair space with \(\nabla e=\eta^{-1}(I-P)\). Put \(V=g_1+e\), \(L_\eta=L_1+\eta^{-1}\), and assume


\[\gamma\le\frac{2}{m+L_\eta}.\]


Let \(\theta^\star\) be a marginal maximizer of \(V\). Let \(Z_0,\xi_0,\xi_1,\dots\) be square-integrable random states with independent standard Gaussian innovations independent of the initialization, let \(\rho_n\) be the law of the \(n\)-th MYIPLA iterate and \(\theta_n\) its parameter block. Then for every \(n\ge0\)


\[\Bigl(\mathbb E\|\theta_n-\theta^\star\|^2\Bigr)^{1/2}\le\sqrt{\frac{d_\theta}{Nm}}+(1-m\gamma)^nW_2(\rho_0,\pi_{N,V})+\frac{33}{20}\,\frac{L_\eta}{m}\sqrt{d_x+\frac{d_\theta}{N}}\,\sqrt\gamma .\]


\textbf{Proof.} \emph{Regularity of \(V\).} By the first step, \(V\) is \(m\)-strongly convex, and \(\nabla V\) is \((L_1+\eta^{-1})=L_\eta\)-Lipschitz; \(V\) is differentiable as a sum of differentiable functions.




\emph{Identification of the laws.} By the previous step (with \(L=L_\eta\)), \(\rho_n\) is the law of the \(n\)-th IPLA iterate for \(V\), with the same initialization and innovations.




\emph{The IPLA bound.} The IPLA estimate with sharp parameter concentration \cite[Section~3.6.2]{LangevinLean2026Proximal} applies to a differentiable \(m\)-strongly convex pair potential with \(L\)-Lipschitz gradient, \(m>0\), step \(0<\gamma\le2/(m+L)\), a marginal maximizer and Gaussian innovations as above. Its hypotheses hold for \(V\) with \(L=L_\eta\), and its conclusion is the displayed inequality. In that estimate the three terms have the following origin: the lifted potential \(F_V\) is \(Nm\)-strongly convex with \(NL_\eta\)-Lipschitz gradient on a space of dimension \(d_\theta+Nd_x\), and the Gaussian Euler step with step \(\gamma/N\) contracts the Wasserstein distance to \(\pi_{N,V}\) by \(1-m\gamma\) per iteration up to a discretization error, which gives the second and third terms; the parameter marginal of \(\pi_{N,V}\) has second moment at most \(d_\theta/(Nm)\) about \(\theta^\star\), which gives the first term after projecting onto the parameter block. \(\square\)




The bound concerns the maximizer \(\theta^\star\) of the \emph{smoothed} marginal likelihood. The smoothing parameter enters through \(L_\eta=L_1+\eta^{-1}\), both in the step restriction and in the discretization term: a smaller \(\eta\) forces a smaller step.

\begin{leanbox}{Lean: \texttt{myipla\_\allowbreak{}parameter\_\allowbreak{}rms\_\allowbreak{}smoothed}}
\leansource{\texttt{myipla\_\allowbreak{}parameter\_\allowbreak{}rms\_\allowbreak{}smoothed} \quad ProximalLangevin/\allowbreak{}MYIPLA.lean, lines 112–143}
\begin{leancode}
\leanline{0}{theorem\ myipla\_parameter\_rms\_smoothed}
\leanline{4}{\{Ω\ :\ Type*\}\ [MeasurableSpace\ Ω]\ (Q\ :\ Measure\ Ω)\ [IsProbabilityMeasure\ Q]}
\leanline{4}{(N\ dθ\ dx\ :\ ℕ)\ (hN\ :\ 0\ <\ N)}
\leanline{4}{\{g₁\ e\ :\ IPLA.Joint\ dθ\ dx\ →\ ℝ\}\ (m\ η\ γ\ :\ ℝ)\ (L₁\ :\ ℝ≥0)}
\leanline{4}{(hg₁\ :\ Differentiable\ ℝ\ g₁)\ (hsc\ :\ StrongConvexOn\ univ\ m\ g₁)}
\leanline{4}{(hLip₁\ :\ LipschitzWith\ L₁\ (gradient\ g₁))}
\leanline{4}{(he\ :\ Differentiable\ ℝ\ e)\ (heconv\ :\ ConvexOn\ ℝ\ univ\ e)}
\leanline{4}{(hLipe\ :\ LipschitzWith\ η.toNNReal⁻¹\ (gradient\ e))}
\leanline{4}{(P\ :\ IPLA.Joint\ dθ\ dx\ →\ IPLA.Joint\ dθ\ dx)}
\leanline{4}{(hgrad\ :\ ∀\ x,\ gradient\ e\ x\ =\ η⁻¹\ •\ (x\ -\ P\ x))}
\leanline{4}{(hm\ :\ 0\ <\ m)\ (hη\ :\ 0\ <\ η)\ (hγ\ :\ 0\ <\ γ)}
\leanline{4}{(hsmall\ :\ γ\ ≤\ 2\ /\ (m\ +\ (smoothingLipschitz\ L₁\ η\ :\ ℝ)))}
\leanline{4}{(θStar\ :\ EuclideanSpace\ ℝ\ (Fin\ dθ))}
\leanline{4}{(hMaximizer\ :\ IPLA.IsMarginalMaximizer\ (fun\ x\ ↦\ g₁\ x\ +\ e\ x)\ θStar)}
\leanline{4}{(initial\ :\ Lp\ (IPLA.State\ N\ dθ\ dx)\ 2\ Q)}
\leanline{4}{(noise\ :\ ℕ\ →\ Lp\ (IPLA.State\ N\ dθ\ dx)\ 2\ Q)}
\leanline{4}{(hGaussian\ :\ IPLA.GaussianInnovations\ Q\ initial\ noise)\ (n\ :\ ℕ)\ :}
\leanline{4}{Real.sqrt\ (∫\ z,\ ‖IPLA.parameterCLM\ z\ -\ θStar‖\ \textasciicircum{}\ 2}
\leanline{6}{∂(myiplaLaws\ Q\ g₁\ P\ η\ γ\ initial\ (fun\ k\ ↦\ noise\ k)\ n))\ ≤}
\leanline{6}{Real.sqrt\ ((dθ\ :\ ℝ)\ /\ ((N\ :\ ℝ)\ *\ m))\ +}
\leanline{6}{(1\ -\ m\ *\ γ)\ \textasciicircum{}\ n\ *\ ULA.wasserstein2Real\ (Q.map\ initial)}
\leanline{8}{(IPLA.gibbsTarget\ (fun\ x\ ↦\ g₁\ x\ +\ e\ x))\ +}
\leanline{6}{(33\ /\ 20\ :\ ℝ)\ *\ ((smoothingLipschitz\ L₁\ η\ :\ ℝ)\ /\ m)\ *}
\leanline{8}{Real.sqrt\ ((dx\ :\ ℝ)\ +\ (dθ\ :\ ℝ)\ /\ (N\ :\ ℝ))\ *\ Real.sqrt\ γ}
\end{leancode}
\leanlabel{How the Lean proof follows the argument}
\texttt{myipla\_\allowbreak{}parameter\_\allowbreak{}rms\_\allowbreak{}smoothed} takes the probability space \texttt{Q} with \texttt{[Is\allowbreak{}Probability\allowbreak{}Measure Q]}, \texttt{h\allowbreak{}N : 0 < N}, the hypotheses \texttt{hg\ensuremath{{}_{1}}}, \texttt{hsc : Strong\allowbreak{}Convex\allowbreak{}On univ m g\ensuremath{{}_{1}}}, \texttt{h\allowbreak{}Lip\ensuremath{{}_{1}} : Lipschitz\allowbreak{}With L\ensuremath{{}_{1}} (gradient g\ensuremath{{}_{1}})}, \texttt{he}, \texttt{heconv : Convex\allowbreak{}On \ensuremath{\mathbb R} univ e}, \texttt{h\allowbreak{}Lipe : Lipschitz\allowbreak{}With \ensuremath{\eta}.\allowbreak{}to\allowbreak{}NNReal\ensuremath{{}^{-1}} (gradient e)} (for \(\eta>0\) this constant is \(\eta^{-1}\)), \texttt{hgrad}, \texttt{hm}, \texttt{h\ensuremath{\eta}}, \texttt{h\ensuremath{\gamma}}, the step restriction \texttt{hsmall : \ensuremath{\gamma} \ensuremath{\le} 2 /\allowbreak{} (m + (smoothing\allowbreak{}Lipschitz L\ensuremath{{}_{1}} \ensuremath{\eta} : \ensuremath{\mathbb R}))}, the maximizer hypothesis \texttt{h\allowbreak{}Maximizer : IPLA.\allowbreak{}Is\allowbreak{}Marginal\allowbreak{}Maximizer (fun x \ensuremath{\mapsto} g\ensuremath{{}_{1}} x + e x) \ensuremath{\theta}\allowbreak{}Star}, and \texttt{h\allowbreak{}Gaussian : IPLA.\allowbreak{}Gaussian\allowbreak{}Innovations Q initial noise}. The last is a structure with three fields: \texttt{coordinate\allowbreak{}Law} (each coordinate of each innovation has law \texttt{gaussian\allowbreak{}Real 0 1}), \texttt{independent} (\texttt{i\allowbreak{}Indep\allowbreak{}Fun}: mutual independence of all coordinates indexed by pairs \((k,a)\)) and \texttt{independent\allowbreak{}Initial} (\texttt{Indep\allowbreak{}Fun}: the initialization is independent of the whole innovation sequence). \texttt{IPLA.\allowbreak{}Is\allowbreak{}Marginal\allowbreak{}Maximizer} compares the extended-real integrals \(\int^-e^{-V(\theta,x)}\,dx\) (lower Lebesgue integrals, \texttt{\ensuremath{\int}\ensuremath{{}^{-}}}, valued in \texttt{\ensuremath{\mathbb R}\ensuremath{\ge}0\ensuremath{\infty}}).




The conclusion bounds \texttt{Real.\allowbreak{}sqrt (\ensuremath{\int} z, \ensuremath{\Vert}IPLA.\allowbreak{}parameter\allowbreak{}CLM z -\allowbreak{} \ensuremath{\theta}\allowbreak{}Star\ensuremath{\Vert} \textasciicircum{} 2 \ensuremath{\partial}(myipla\allowbreak{}Laws …))}, i.e. \((\mathbb E\|\theta_n-\theta^\star\|^2)^{1/2}\) written as an integral against \(\rho_n\). The transport term is \texttt{ULA.\allowbreak{}wasserstein2\allowbreak{}Real (Q.\allowbreak{}map initial) (IPLA.\allowbreak{}gibbs\allowbreak{}Target (fun x \ensuremath{\mapsto} g\ensuremath{{}_{1}} x + e x))}, the real number obtained from the extended-real Wasserstein distance; that distance is finite here (square-integrable initialization, target with finite second moment), as is established inside the cited IPLA estimate.




The proof has three lines: \texttt{h\allowbreak{}Lip} is \texttt{smoothed\_\allowbreak{}gradient\_\allowbreak{}lipschitz hg\ensuremath{{}_{1}} he h\allowbreak{}Lip\ensuremath{{}_{1}} h\allowbreak{}Lipe}, the law is rewritten with \texttt{myipla\allowbreak{}Laws\_\allowbreak{}eq\_\allowbreak{}ipla\allowbreak{}Laws}, and the result is \texttt{sharp\_\allowbreak{}ipla\_\allowbreak{}parameter\_\allowbreak{}rms} applied to \(V\) with constant \texttt{smoothing\allowbreak{}Lipschitz L\ensuremath{{}_{1}} \ensuremath{\eta}} and strong convexity \texttt{smoothed\_\allowbreak{}strong\allowbreak{}Convex hsc heconv}.
\end{leanbox}


Earlier results used here: \cite[Section~3.17.4]{LangevinLean2026ULASC}; \cite[Section~3.1.1]{LangevinLean2026IPLA}; \cite[Section~3.1.6]{LangevinLean2026IPLA}; \cite[Section~3.3.1]{LangevinLean2026IPLA}; \cite[Section~3.3.2]{LangevinLean2026IPLA}.

\section{\texorpdfstring{The theorem for the Moreau envelope}{The theorem for the Moreau envelope}}


\textbf{Claim.} Let \(N\ge1\), \(\eta>0\), \(L_1\ge0\), let \(g_1\) be differentiable with \(L_1\)-Lipschitz gradient, and let \(g_2\) be convex and bounded below on the (finite-dimensional) pair space. Let \(P_\eta\) be the proximal map of \(g_2\), \(V_\eta=g_1+g_2^\eta\), and let \(\rho_n\) be the law of the \(n\)-th MYIPLA iterate with proximal map \(P_\eta\), smoothing parameter \(\eta\), any step \(\gamma\in\mathbb R\), and square-integrable initialization and innovations. Then for every \(n\):



\begin{enumerate}


\item \(\rho_n\) is a probability measure;

\item \(\int\|\theta(z)-\vartheta\|^2\,\rho_n(dz)<\infty\) for every \(\vartheta\in\mathbb R^{d_\theta}\);

\item if in addition \(g_1\) is \(m\)-strongly convex with \(m>0\), \(0<\gamma\le2/(m+L_\eta)\) with \(L_\eta=L_1+\eta^{-1}\), \(\theta^\star_\eta\) is a marginal maximizer of \(V_\eta\), and the innovations are independent standard Gaussian and independent of the initialization, then
\[\Bigl(\mathbb E\|\theta_n-\theta^\star_\eta\|^2\Bigr)^{1/2}\le\sqrt{\frac{d_\theta}{Nm}}+(1-m\gamma)^nW_2(\rho_0,\pi_{N,V_\eta})+\frac{33}{20}\,\frac{L_\eta}{m}\sqrt{d_x+\frac{d_\theta}N}\,\sqrt\gamma .\]

\end{enumerate}



\textbf{Proof.} By the Moreau envelope section, under the stated hypotheses on \(g_2\) and \(\eta>0\): \(e=g_2^\eta\) is differentiable and convex, \(\nabla e=\eta^{-1}(I-P_\eta)\), and \(\nabla e\) is \(\eta^{-1}\)-Lipschitz (the residual is \(\eta^{-1}\)-Lipschitz). With \(P=P_\eta\) all hypotheses of the previous two steps that concern \(e\) and \(P\) are therefore satisfied, and \(\nabla(g_1+e)\) is \(L_\eta\)-Lipschitz by the first step.




(a) By the second step, \(\rho_n\) is the pushforward of the probability measure \(Q\) under a measurable map \(Y_n\), so \(\rho_n(\text{whole space})=Q(Y_n^{-1}(\text{whole space}))=Q(\Omega)=1\).




(b) is part (b) of the second step.




(c) is the previous step with \(e=g_2^\eta\) and \(P=P_\eta\). \(\square\)




All properties of the envelope used here are consequences of convexity and the lower bound of \(g_2\); none is assumed. Parts (a) and (b) need neither strong convexity, nor Gaussian innovations, nor a step restriction.

\begin{leanbox}{Lean: \texttt{moreau\_\allowbreak{}myipla\_\allowbreak{}parameter\_\allowbreak{}rms}, \texttt{moreau\_\allowbreak{}myipla\_\allowbreak{}parameter\_\allowbreak{}memLp}, \texttt{moreau\_\allowbreak{}myipla\_\allowbreak{}isProbability}}
\leansource{\texttt{moreau\_\allowbreak{}myipla\_\allowbreak{}parameter\_\allowbreak{}rms} \quad ProximalLangevin/\allowbreak{}MYIPLA.lean, lines 150–181}
\begin{leancode}
\leanline{0}{theorem\ moreau\_myipla\_parameter\_rms}
\leanline{4}{\{Ω\ :\ Type*\}\ [MeasurableSpace\ Ω]\ (Q\ :\ Measure\ Ω)\ [IsProbabilityMeasure\ Q]}
\leanline{4}{(N\ dθ\ dx\ :\ ℕ)\ (hN\ :\ 0\ <\ N)}
\leanline{4}{\{g₁\ g₂\ :\ IPLA.Joint\ dθ\ dx\ →\ ℝ\}\ (m\ η\ γ\ :\ ℝ)\ (L₁\ :\ ℝ≥0)}
\leanline{4}{(hg₁\ :\ Differentiable\ ℝ\ g₁)\ (hsc\ :\ StrongConvexOn\ univ\ m\ g₁)}
\leanline{4}{(hLip₁\ :\ LipschitzWith\ L₁\ (gradient\ g₁))}
\leanline{4}{(hg₂\ :\ ConvexOn\ ℝ\ univ\ g₂)\ (hb\ :\ ∃\ b\ :\ ℝ,\ ∀\ z,\ b\ ≤\ g₂\ z)}
\leanline{4}{(hm\ :\ 0\ <\ m)\ (hη\ :\ 0\ <\ η)\ (hγ\ :\ 0\ <\ γ)}
\leanline{4}{(hsmall\ :\ γ\ ≤\ 2\ /\ (m\ +\ (smoothingLipschitz\ L₁\ η\ :\ ℝ)))}
\leanline{4}{(θStar\ :\ EuclideanSpace\ ℝ\ (Fin\ dθ))}
\leanline{4}{(hMaximizer\ :\ IPLA.IsMarginalMaximizer\ (smoothedPotential\ g₁\ g₂\ η)\ θStar)}
\leanline{4}{(initial\ :\ Lp\ (IPLA.State\ N\ dθ\ dx)\ 2\ Q)}
\leanline{4}{(noise\ :\ ℕ\ →\ Lp\ (IPLA.State\ N\ dθ\ dx)\ 2\ Q)}
\leanline{4}{(hGaussian\ :\ IPLA.GaussianInnovations\ Q\ initial\ noise)\ (n\ :\ ℕ)\ :}
\leanline{4}{Real.sqrt\ (∫\ z,\ ‖IPLA.parameterCLM\ z\ -\ θStar‖\ \textasciicircum{}\ 2}
\leanline{6}{∂(myiplaLaws\ Q\ g₁\ (prox\ g₂\ η)\ η\ γ\ initial\ (fun\ k\ ↦\ noise\ k)\ n))\ ≤}
\leanline{6}{Real.sqrt\ ((dθ\ :\ ℝ)\ /\ ((N\ :\ ℝ)\ *\ m))\ +}
\leanline{6}{(1\ -\ m\ *\ γ)\ \textasciicircum{}\ n\ *\ ULA.wasserstein2Real\ (Q.map\ initial)}
\leanline{8}{(IPLA.gibbsTarget\ (smoothedPotential\ g₁\ g₂\ η))\ +}
\leanline{6}{(33\ /\ 20\ :\ ℝ)\ *\ ((smoothingLipschitz\ L₁\ η\ :\ ℝ)\ /\ m)\ *}
\leanline{8}{Real.sqrt\ ((dx\ :\ ℝ)\ +\ (dθ\ :\ ℝ)\ /\ (N\ :\ ℝ))\ *\ Real.sqrt\ γ}
\end{leancode}
\leansource{\texttt{moreau\_\allowbreak{}myipla\_\allowbreak{}parameter\_\allowbreak{}memLp} \quad ProximalLangevin/\allowbreak{}MYIPLA.lean, lines 184–203}
\begin{leancode}
\leanline{0}{theorem\ moreau\_myipla\_parameter\_memLp}
\leanline{4}{\{Ω\ :\ Type*\}\ [MeasurableSpace\ Ω]\ (Q\ :\ Measure\ Ω)\ [IsProbabilityMeasure\ Q]}
\leanline{4}{\{N\ dθ\ dx\ :\ ℕ\}\ (hN\ :\ 0\ <\ N)}
\leanline{4}{\{g₁\ g₂\ :\ IPLA.Joint\ dθ\ dx\ →\ ℝ\}\ (η\ γ\ :\ ℝ)\ (L₁\ :\ ℝ≥0)}
\leanline{4}{(hg₁\ :\ Differentiable\ ℝ\ g₁)\ (hLip₁\ :\ LipschitzWith\ L₁\ (gradient\ g₁))}
\leanline{4}{(hg₂\ :\ ConvexOn\ ℝ\ univ\ g₂)\ (hb\ :\ ∃\ b\ :\ ℝ,\ ∀\ z,\ b\ ≤\ g₂\ z)\ (hη\ :\ 0\ <\ η)}
\leanline{4}{(initial\ :\ Lp\ (IPLA.State\ N\ dθ\ dx)\ 2\ Q)}
\leanline{4}{(noise\ :\ ℕ\ →\ Lp\ (IPLA.State\ N\ dθ\ dx)\ 2\ Q)\ (n\ :\ ℕ)}
\leanline{4}{(θ\ :\ EuclideanSpace\ ℝ\ (Fin\ dθ))\ :}
\leanline{4}{MemLp\ (fun\ z\ ↦\ IPLA.parameterCLM\ z\ -\ θ)\ 2}
\leanline{6}{(myiplaLaws\ Q\ g₁\ (prox\ g₂\ η)\ η\ γ\ initial\ (fun\ k\ ↦\ noise\ k)\ n)}
\end{leancode}
\leansource{\texttt{moreau\_\allowbreak{}myipla\_\allowbreak{}isProbability} \quad ProximalLangevin/\allowbreak{}MYIPLA.lean, lines 207–228}
\begin{leancode}
\leanline{0}{theorem\ moreau\_myipla\_isProbability}
\leanline{4}{\{Ω\ :\ Type*\}\ [MeasurableSpace\ Ω]\ (Q\ :\ Measure\ Ω)\ [IsProbabilityMeasure\ Q]}
\leanline{4}{\{N\ dθ\ dx\ :\ ℕ\}\ (hN\ :\ 0\ <\ N)}
\leanline{4}{\{g₁\ g₂\ :\ IPLA.Joint\ dθ\ dx\ →\ ℝ\}\ (η\ γ\ :\ ℝ)\ (L₁\ :\ ℝ≥0)}
\leanline{4}{(hg₁\ :\ Differentiable\ ℝ\ g₁)\ (hLip₁\ :\ LipschitzWith\ L₁\ (gradient\ g₁))}
\leanline{4}{(hg₂\ :\ ConvexOn\ ℝ\ univ\ g₂)\ (hb\ :\ ∃\ b\ :\ ℝ,\ ∀\ z,\ b\ ≤\ g₂\ z)\ (hη\ :\ 0\ <\ η)}
\leanline{4}{(initial\ :\ Lp\ (IPLA.State\ N\ dθ\ dx)\ 2\ Q)}
\leanline{4}{(noise\ :\ ℕ\ →\ Lp\ (IPLA.State\ N\ dθ\ dx)\ 2\ Q)\ (n\ :\ ℕ)\ :}
\leanline{4}{IsProbabilityMeasure\ (myiplaLaws\ Q\ g₁\ (prox\ g₂\ η)\ η\ γ\ initial\ (fun\ k\ ↦\ noise\ k)\ n)}
\end{leancode}
\leanlabel{How the Lean proof follows the argument}
In all three theorems the envelope facts are produced in the same way: \texttt{gradient\_\allowbreak{}moreau hg\ensuremath{{}_{2}} hb h\ensuremath{\eta}} gives \texttt{hgrad}, \(\nabla g_2^\eta(x)=\eta^{-1}(x-P_\eta(x))\); \texttt{hgeq} upgrades it to the equality of functions \(\nabla g_2^\eta=r_\eta\) by \texttt{funext}; and \texttt{h\allowbreak{}Lipe}, the \(\eta^{-1}\)-Lipschitz property of \(\nabla g_2^\eta\), is \texttt{residual\_\allowbreak{}lipschitz hg\ensuremath{{}_{2}} hb h\ensuremath{\eta}} after rewriting with \texttt{hgeq} and converting \((\eta^{-1})^+\) into \((\eta^+)^{-1}\) with \texttt{Real.\allowbreak{}to\allowbreak{}NNReal\_\allowbreak{}inv}. Differentiability and convexity of the envelope are \texttt{moreau\_\allowbreak{}differentiable} and \texttt{moreau\_\allowbreak{}convex}. The proximal map is \texttt{prox g\ensuremath{{}_{2}} \ensuremath{\eta}} and the smoothed potential \texttt{smoothed\allowbreak{}Potential g\ensuremath{{}_{1}} g\ensuremath{{}_{2}} \ensuremath{\eta}}.



\begin{enumerate}


\item \texttt{moreau\_\allowbreak{}myipla\_\allowbreak{}parameter\_\allowbreak{}rms} is (c): it applies \texttt{myipla\_\allowbreak{}parameter\_\allowbreak{}rms\_\allowbreak{}smoothed} with these facts. Its hypothesis list contains only \texttt{hg\ensuremath{{}_{2}} : Convex\allowbreak{}On \ensuremath{\mathbb R} univ g\ensuremath{{}_{2}}} and \texttt{hb} for the nonsmooth part.

\item \texttt{moreau\_\allowbreak{}myipla\_\allowbreak{}parameter\_\allowbreak{}mem\allowbreak{}Lp} is (b): it applies \texttt{myipla\allowbreak{}Laws\_\allowbreak{}parameter\_\allowbreak{}mem\allowbreak{}Lp} with the Lipschitz constant \texttt{smoothing\allowbreak{}Lipschitz L\ensuremath{{}_{1}} \ensuremath{\eta}} supplied by \texttt{smoothed\_\allowbreak{}gradient\_\allowbreak{}lipschitz}; the step \texttt{\ensuremath{\gamma}} is an arbitrary real number.

\item \texttt{moreau\_\allowbreak{}myipla\_\allowbreak{}is\allowbreak{}Probability} is (a). After rewriting the law with \texttt{myipla\allowbreak{}Laws\_\allowbreak{}eq\_\allowbreak{}ipla\allowbreak{}Laws} and unfolding \texttt{IPLA.\allowbreak{}laws}, the goal \texttt{Is\allowbreak{}Probability\allowbreak{}Measure (Q.\allowbreak{}map …)} is proved by \texttt{constructor} and the computation \texttt{Measure.\allowbreak{}map\_\allowbreak{}apply\_\allowbreak{}of\_\allowbreak{}aemeasurable} (evaluation of a pushforward on a measurable set, here \texttt{Measurable\allowbreak{}Set.\allowbreak{}univ}), \texttt{Set.\allowbreak{}preimage\_\allowbreak{}univ} and \texttt{measure\_\allowbreak{}univ}; measurability of the \(L^2\) iterate is \texttt{(Lp.\allowbreak{}aestrongly\allowbreak{}Measurable \_\allowbreak{}).\allowbreak{}aemeasurable}.

\end{enumerate}
\end{leanbox}


Earlier results used here: \cite[Section~3.17.4]{LangevinLean2026ULASC}; \cite[Section~3.1.1]{LangevinLean2026IPLA}; \cite[Section~3.1.6]{LangevinLean2026IPLA}; \cite[Section~3.3.1]{LangevinLean2026IPLA}; \cite[Section~3.3.2]{LangevinLean2026IPLA}; \cite[Section~3.3.3]{LangevinLean2026IPLA}.

\begin{thebibliography}{ULA-LSI}
\bibitem[ULA-SC]{LangevinLean2026ULASC}
LangevinLean Project. \emph{Unadjusted Langevin under Strong Convexity}. LL-01, 2026. Draft report, content revision d34ead06987c. \url{https://odakyildiz.github.io/LangevinLean/docs/pdf/UnadjustedLangevin/FullProof.pdf}.
\bibitem[IPLA-SC]{LangevinLean2026IPLA}
LangevinLean Project. \emph{Interacting Particle Langevin under Strong Convexity}. LL-03, 2026. Draft report, content revision ec0201df2207. \url{https://odakyildiz.github.io/LangevinLean/docs/pdf/IPLA/FullProof.pdf}.
\bibitem[Proximal]{LangevinLean2026Proximal}
LangevinLean Project. \emph{Proximal particle Langevin}. LL-05, 2026. Draft report, content revision b5e97fa352ba. \url{https://odakyildiz.github.io/LangevinLean/docs/pdf/ProximalLangevin/FullProof.pdf}.
\end{thebibliography}

\end{document}
